\backmatter
\AfterBibliographyPreamble{%
Every measurement in this work is built on somebody's published result, and the
first purpose of this chapter is that those people are credited by name. A surname
in a footnote is not a citation. The second purpose is that the credit is checkable:
each source is registered in \texttt{anchor\_sift\_citations}, where the copy sits
itself, and this bibliography is generated from that registry and pinned to one
commit of it. A reference therefore resolves not to a title but to exact bytes.

Credit here does not depend on agreement. Where this work reproduces a published
result it says so, and where it fails to reproduce one it says that too, under the
same names, because a refutation rests on the original as much as a confirmation does.

\begin{description}
\item[Registry commit] \texttt{2aae1103d28d9b919eae59cc9062aeaffe1fa0b2}
\item[Dated] 2026-10-01
\item[Inventory] \texttt{MANIFEST.tsv}, signed; the hashes below are its rows
\end{description}

A reader who disagrees with a number in this research paper can ask which copy it was checked
against, and the hash answers exactly. Two copies of a paper are rarely the same
bytes.
\par\medskip\raggedright}
\KOMAoptions{bibliography=totoc}
\begin{thebibliography}{61}

\bibitem{src:Aaronson-2022}
S. Aaronson. \emph{On black holes, holography, the Quantum Extended Church-Turing Thesis, fully homomorphic encryption, and brain uploading}. Shtetl-Optimized, 27 July 2022, \texttt{https:/\allowbreak{}/\allowbreak{}scottaaronson.blog/\allowbreak{}?p=6599}.

\bibitem{src:Almheiri-Dong-and-Harlow-2015}
A. Almheiri, X. Dong, D. Harlow. \emph{Bulk locality and quantum error correction in AdS/CFT}. JHEP 04, 163, \texttt{doi:10.1007/\allowbreak{}jhep04(2015)163}.

\bibitem{src:Arnold-1957}
V. I. Arnold. \emph{On functions of three variables}. Doklady Akademii Nauk SSSR 114(4) 679-681. 1957.

\bibitem{src:Barsi-and-Maestrini-1973}
F. Barsi, P. Maestrini. \emph{Error correcting properties of redundant residue number systems}. IEEE Transactions on Computers C-22(3), 307-315, \texttt{doi:10.1109/\allowbreak{}t-c.1973.223711}.

\bibitem{src:Bennett-1973}
C. H. Bennett. \emph{Logical reversibility of computation}. IBM Journal of Research and Development 17(6) 525-532, \texttt{doi:10.1147/\allowbreak{}rd.176.0525}. 1973.

\bibitem{src:Borsten-and-Kim-2026}
L. Borsten, H. Kim. \emph{Limits to Computational Acceleration Imposed by Quantum Field Theory and Quantum Gravity}. arXiv:2604.00182. 2026.

\bibitem{src:Borwein}
J. Borwein, W. Galway, D. Borwein. \emph{Finding and Excluding b-ary Machin-Type Individual Digit Formulae}. Canadian Journal of Mathematics 56(5) 897-925, \texttt{doi:10.4153/\allowbreak{}cjm-2004-041-2}.

\bibitem{src:Chaitin-1975}
G. J. Chaitin. \emph{A theory of program size formally identical to information theory}. Journal of the ACM 22, \texttt{doi:10.1145/\allowbreak{}321892.321894}. 1975.

\bibitem{src:Chaitin-1987}
G. J. Chaitin. \emph{Algorithmic Information Theory}. Cambridge University Press. 1987.

\bibitem{src:Chang-1959}
C. C. Chang. \emph{On unions of chains of models}. Proceedings of the American Mathematical Society 10, \texttt{doi:10.1090/\allowbreak{}s0002-9939-1959-0103812-3}.

\bibitem{src:Fefferman}
Charles L. Fefferman. \emph{Existence and Smoothness of the Navier-Stokes Equation}. Clay Mathematics Institute, Millennium Prize Problem statement, \texttt{https:/\allowbreak{}/\allowbreak{}www.claymath.org/\allowbreak{}wp-content/\allowbreak{}uploads/\allowbreak{}2022/\allowbreak{}06/\allowbreak{}navierstokes.pdf}.\newline
{\footnotesize Copy held \texttt{sources/\allowbreak{}mathematics/\allowbreak{}millennium\_\allowbreak{}problems/\allowbreak{}clay\_\allowbreak{}fefferman\_\allowbreak{}navier\_\allowbreak{}stokes.pdf}, SHA-256 \texttt{c1b5f27b1a64705cfaf1afceea513db5deedca8a18ca56ab32e7f86445a06d0c}.}

\bibitem{src:Fisher}
R. A. Fisher, F. Yates. \emph{Statistical Tables for Biological, Agricultural and Medical Research}. Oliver and Boyd, Edinburgh. 1938.

\bibitem{src:Garner-1959}
H. L. Garner. \emph{The residue number system}. IRE Transactions on Electronic Computers EC-8, 140-147, \texttt{doi:10.1109/\allowbreak{}tec.1959.5219515}.

\bibitem{src:Gentzen-1936}
G. Gentzen. \emph{Die Widerspruchsfreiheit der reinen Zahlentheorie}. Mathematische Annalen 112, \texttt{doi:10.1007/\allowbreak{}bf01565428}. 1936.

\bibitem{src:Goodstein-1944}
R. L. Goodstein. \emph{On the restricted ordinal theorem}. Journal of Symbolic Logic 9, \texttt{doi:10.2307/\allowbreak{}2268019}. 1944.

\bibitem{src:Gouv-a-1997}
F. Q. Gouvêa. \emph{p-adic Numbers: An Introduction, second edition}. Springer. 1997.

\bibitem{src:Hamkins-and-Lewis-2000}
J. D. Hamkins, A. Lewis. \emph{Infinite time Turing machines}. Journal of Symbolic Logic 65, \texttt{doi:10.2307/\allowbreak{}2586556}. 2000.

\bibitem{src:Hewitt-and-Ross-1963}
E. Hewitt, K. A. Ross. \emph{Abstract Harmonic Analysis I}. Springer. 1963.

\bibitem{src:Hodges-1993}
W. Hodges. \emph{Model Theory}. Cambridge University Press. 1993.

\bibitem{src:Hope-1968}
A. C. A. Hope. \emph{A simplified Monte Carlo significance test procedure}. Journal of the Royal Statistical Society, Series B 30, \texttt{doi:10.1111/\allowbreak{}j.2517-6161.1968.tb00759.x}.

\bibitem{src:Hurwitz-1891}
A. Hurwitz. \emph{Ueber die angenäherte Darstellung der Irrationalzahlen durch rationale Brüche}. Mathematische Annalen 39, \texttt{doi:10.1007/\allowbreak{}bf01206656}. 1891.

\bibitem{src:Jones-1706}
W. Jones. \emph{Synopsis Palmariorum Matheseos}. London. 1706.

\bibitem{src:Kac}
M. Kac. \emph{On the notion of recurrence in discrete stochastic processes}. Bulletin of the American Mathematical Society 53(10) 1002-1010, \texttt{doi:10.1090/\allowbreak{}s0002-9904-1947-08927-8}.

\bibitem{src:Khinchin-1964}
A. Ya. Khinchin. \emph{Continued Fractions}. University of Chicago Press. 1964.

\bibitem{src:Kirby-and-Paris-1982}
L. Kirby, J. Paris. \emph{Accessible independence results for Peano arithmetic}. Bulletin of the London Mathematical Society 14, \texttt{doi:10.1112/\allowbreak{}blms/\allowbreak{}14.4.285}. 1982.

\bibitem{src:Knuth}
D. E. Knuth.

\bibitem{src:Koblitz-1984}
N. Koblitz. \emph{p-adic Numbers, p-adic Analysis, and Zeta-Functions, second edition}. Springer. 1984.

\bibitem{src:Koepke-2005}
P. Koepke. \emph{Turing computations on ordinals}. Bulletin of Symbolic Logic 11, \texttt{doi:10.2178/\allowbreak{}bsl/\allowbreak{}1122038993}. 2005.

\bibitem{src:Kolmogorov}
A. N. Kolmogorov. \emph{A new metric invariant of transitive dynamical systems and automorphisms of Lebesgue spaces}. Dokl. Akad. Nauk SSSR 119(5):861-864, Math-Net.Ru dan22922. 1958.\newline
{\footnotesize Copy held \texttt{sources/\allowbreak{}mathematics/\allowbreak{}ergodic\_\allowbreak{}theory/\allowbreak{}dan22922\_\allowbreak{}kolmogorov\_\allowbreak{}1958.pdf}, SHA-256 \texttt{abf376fa2e2aefaf1492308a8808d79be3431dae4c821ca546719f35a1d4bf85}.}

\bibitem{src:Kolmogorov-1957}
A. N. Kolmogorov. \emph{On the representation of continuous functions of many variables by superposition of continuous functions of one variable and addition}. Doklady Akademii Nauk SSSR 114(5) 953-956. 1957.

\bibitem{src:Kolmogorov-1965}
A. N. Kolmogorov. \emph{Three approaches to the quantitative definition of information}. Problemy Peredachi Informatsii 1(1) 3-11; in English, International Journal of Computer Mathematics 2(1-4) 157-168, 1968, \texttt{doi:10.1080/\allowbreak{}00207166808803030}. 1965.

\bibitem{src:Lambert-1761}
J. H. Lambert. \emph{Mémoire sur quelques propriétés remarquables des quantités transcendantes circulaires et logarithmiques}. Mémoires de l'Académie royale des sciences de Berlin 17, 1761, published 1768.

\bibitem{src:Lancret-1806}
M. A. Lancret. \emph{Mémoire sur les courbes à double courbure}. Mémoires présentés à l'Institut 1. 1806.

\bibitem{src:L-vy-1936}
P. Lévy. \emph{Sur le développement en fraction continue d'un nombre choisi au hasard}. Compositio Mathematica 3. 1936.

\bibitem{src:Lindemann-1882}
F. Lindemann. \emph{Über die Zahl π}. Mathematische Annalen 20, \texttt{doi:10.1007/\allowbreak{}bf01446522}. 1882.

\bibitem{src:Liu-et-al-2024}
Z. Liu, Y. Wang, S. Vaidya, F. Ruehle, J. Halverson, M. Soljačić, T. Y. Hou, M. Tegmark. \emph{KAN: Kolmogorov-Arnold Networks}. arXiv:2404.19756. 2024.

\bibitem{src:o-and-Suszko-1957}
J. Łoś, R. Suszko. \emph{On the extending of models (IV)}. Fundamenta Mathematicae 44, \texttt{doi:10.4064/\allowbreak{}fm-44-1-52-60}. 1957.

\bibitem{src:Martin-L-f-1966}
P. Martin-Löf. \emph{The definition of random sequences}. Information and Control 9, \texttt{doi:10.1016/\allowbreak{}s0019-9958(66)80018-9}. 1966.

\bibitem{src:Ostrowski-1916}
A. Ostrowski. \emph{Über einige Lösungen der Funktionalgleichung φ(x)·φ(y) = φ(xy)}. Acta Mathematica 41, \texttt{doi:10.1007/\allowbreak{}bf02422947}. 1916.

\bibitem{src:Pastawski-Yoshida-Harlow-and-Preskill-2015}
F. Pastawski, B. Yoshida, D. Harlow, J. Preskill. \emph{Holographic quantum error-correcting codes: toy models for the bulk/boundary correspondence}. JHEP 06, 149, \texttt{doi:10.1007/\allowbreak{}jhep06(2015)149}.

\bibitem{src:Peterson-and-Weldon-1972}
W. W. Peterson, E. J. Weldon. \emph{Error-Correcting Codes, second edition}. MIT Press. 1972.

\bibitem{src:Post}
E. L. Post. \emph{Formal Reductions of the General Combinatorial Decision Problem}. American Journal of Mathematics 65(2) 197-215. 1943.

\bibitem{src:Rauzy-1979}
G. Rauzy. \emph{Échanges d'intervalles et transformations induites}. Acta Arithmetica 34, \texttt{doi:10.4064/\allowbreak{}aa-34-4-315-328}. 1979.

\bibitem{src:Reiger-1960}
S. Reiger. \emph{Codes for the correction of 'clustered' errors}. IRE Transactions on Information Theory 6, \texttt{doi:10.1109/\allowbreak{}tit.1960.1057558}.

\bibitem{src:Schoenberg-1946}
I. J. Schoenberg. \emph{Contributions to the problem of approximation of equidistant data by analytic functions. Part A}. Quarterly of Applied Mathematics 4(1) 45-99, \texttt{doi:10.1090/\allowbreak{}qam/\allowbreak{}15914}. 1946.

\bibitem{src:Shannon}
C. E. Shannon. \emph{A Mathematical Theory of Communication}. Bell System Technical Journal 27, 379-423 and 623-656, July and October 1948.\newline
{\footnotesize Copy held \texttt{sources/\allowbreak{}physics/\allowbreak{}information\_\allowbreak{}theory/\allowbreak{}shannon\_\allowbreak{}1948\_\allowbreak{}mathematical\_\allowbreak{}theory\_\allowbreak{}of\_\allowbreak{}communication.pdf}, SHA-256 \texttt{6e4e3411984f3edf99dbfe8b941cb5e8a321379ff0cae6ae5c1f592ad8882ca8}.}

\bibitem{src:Shoenfield-1959}
J. R. Shoenfield. \emph{On degrees of unsolvability}. Annals of Mathematics 69, \texttt{doi:10.2307/\allowbreak{}1970028}. 1959.

\bibitem{src:Slater-1950}
N. B. Slater. \emph{The distribution of the integers N for which \{θN\} < φ}. Proceedings of the Cambridge Philosophical Society 46, \texttt{doi:10.1017/\allowbreak{}s0305004100026086}. 1950.

\bibitem{src:Slater-1967}
N. B. Slater. \emph{Gaps and steps for the sequence nθ mod 1}. Proceedings of the Cambridge Philosophical Society 63, 1115-1123, \texttt{doi:10.1017/\allowbreak{}s0305004100042195}. 1967.

\bibitem{src:Sorkin-1994}
R. D. Sorkin. \emph{Quantum mechanics as quantum measure theory}. Modern Physics Letters A 9, \texttt{doi:10.1142/\allowbreak{}s021773239400294x}. 1994.

\bibitem{src:S-s-1958}
V. T. Sós. \emph{On the distribution mod 1 of the sequence nα}. Annales Universitatis Scientiarum Budapestinensis de Rolando Eötvös Nominatae, Sectio Mathematica 1. 1958.

\bibitem{src:Sweldens-1996}
W. Sweldens. \emph{The lifting scheme: a custom-design construction of biorthogonal wavelets}. Applied and Computational Harmonic Analysis 3, \texttt{doi:10.1006/\allowbreak{}acha.1996.0015}.

\bibitem{src:Szabo-and-Tanaka-1967}
N. S. Szabo, R. I. Tanaka. \emph{Residue Arithmetic and Its Applications to Computer Technology}. McGraw-Hill. 1967.

\bibitem{src:Tate-1950}
J. Tate. \emph{Fourier analysis in number fields and Hecke's zeta-functions}. thesis, Princeton, 1950; in J. W. S. Cassels and A. Fröhlich (eds.), Algebraic Number Theory, Academic Press, 1967.

\bibitem{src:van-Dantzig-1930}
D. van Dantzig. \emph{Über topologisch homogene Kontinua}. Fundamenta Mathematicae 15. 1930.

\bibitem{src:van-Ravenstein-1988}
T. van Ravenstein. \emph{The three gap theorem (Steinhaus conjecture)}. Journal of the Australian Mathematical Society, Series A 45, \texttt{doi:10.1017/\allowbreak{}s1446788700031062}. 1988.

\bibitem{src:Vietoris-1927}
L. Vietoris. \emph{Über den höheren Zusammenhang kompakter Räume und eine Klasse von zusammenhangstreuen Abbildungen}. Mathematische Annalen 97. 1927.

\bibitem{src:Walters-1982}
P. Walters. \emph{An Introduction to Ergodic Theory}. Springer. 1982.

\bibitem{src:Watson-and-Hastings-1966}
R. W. Watson, C. W. Hastings. \emph{Self-checked computation using residue arithmetic}. Proceedings of the IEEE 54(12), 1920-1931, \texttt{doi:10.1109/\allowbreak{}proc.1966.5275}.

\bibitem{src:Weyl-1916}
H. Weyl. \emph{Über die Gleichverteilung von Zahlen mod. Eins}. Mathematische Annalen 77, \texttt{doi:10.1007/\allowbreak{}bf01475864}. 1916.

\bibitem{src:Wootters-and-Zurek-1982}
W. K. Wootters, W. H. Zurek. \emph{A single quantum cannot be cloned}. Nature 299, \texttt{doi:10.1038/\allowbreak{}299802a0}. 1982.

\end{thebibliography}
